%% Response to the reviewers, revision 1 of CAOR-D-26-01885.
%% Build (after compiling ../main.tex, whose .aux supplies the cross
%% references): pdflatex response_letter (twice)
\documentclass[11pt]{article}
\usepackage[a4paper,margin=2.4cm]{geometry}
\usepackage{amsmath,amssymb}
\usepackage{xcolor}
\usepackage{parskip}
\usepackage{enumitem}
\usepackage{xr}
\usepackage[colorlinks=true,allcolors=blue!50!black]{hyperref}
\externaldocument[m-]{../main}
% generated by make_tables.py — do not hand-edit
\newcommand{\batonGrandBest}{53.7\%}
\newcommand{\budCityDp}{12.2\%}
\newcommand{\budCityDpNLarge}{11.9\%}
\newcommand{\budCityDpNMid}{7.0\%}
\newcommand{\budCityDpNSmall}{-0.3\%}
\newcommand{\budCityDpThree}{12.3\%}
\newcommand{\budCityLarge}{12.9\%}
\newcommand{\budCityMid}{10.5\%}
\newcommand{\budCitySmall}{3.8\%}
\newcommand{\budCityThrLarge}{12.0\%}
\newcommand{\budCityThrMid}{10.0\%}
\newcommand{\budCityThrSmall}{0.4\%}
\newcommand{\budDethDetDp}{26.4\%}
\newcommand{\budDethDetDpNLarge}{26.3\%}
\newcommand{\budDethDetDpNMid}{25.7\%}
\newcommand{\budDethDetDpNSmall}{22.7\%}
\newcommand{\budDethDetDpThree}{29.5\%}
\newcommand{\budDethDetLarge}{28.3\%}
\newcommand{\budDethDetMid}{27.9\%}
\newcommand{\budDethDetSmall}{25.8\%}
\newcommand{\budDethDetThrLarge}{25.0\%}
\newcommand{\budDethDetThrMid}{24.8\%}
\newcommand{\budDethDetThrSmall}{23.0\%}
\newcommand{\budDethSaaDp}{22.2\%}
\newcommand{\budDethSaaDpNLarge}{22.0\%}
\newcommand{\budDethSaaDpNMid}{12.1\%}
\newcommand{\budDethSaaDpNSmall}{-1.6\%}
\newcommand{\budDethSaaDpThree}{56.4\%}
\newcommand{\budDethSaaLarge}{56.3\%}
\newcommand{\budDethSaaMid}{53.6\%}
\newcommand{\budDethSaaSmall}{41.7\%}
\newcommand{\budDethSaaThrLarge}{19.6\%}
\newcommand{\budDethSaaThrMid}{16.5\%}
\newcommand{\budDethSaaThrSmall}{-2.4\%}
\newcommand{\budgetCross}{100}
\newcommand{\certBeatN}{17/40}
\newcommand{\certGap}{9.34\%}
\newcommand{\certGapHighs}{11.03\%}
\newcommand{\certGapMax}{12.55\%}
\newcommand{\cfGapHi}{1.3}
\newcommand{\cfGapLo}{0.1}
\newcommand{\cityActShare}{37\%}
\newcommand{\cityDetK}{8.5}
\newcommand{\cityDetReact}{36.8}
\newcommand{\citySaaK}{12.1}
\newcommand{\citySaaPem}{0.05\%}
\newcommand{\citySaaReact}{0.04}
\newcommand{\ctdBaton}{51.9\%}
\newcommand{\ctdBatonFull}{69.1\%}
\newcommand{\ctdVOne}{0.0\%}
\newcommand{\dearSbBaton}{42.6\%}
\newcommand{\dearSbOrc}{24.8\%}
\newcommand{\dearSbThresh}{8.8\%}
\newcommand{\depCfAbove}{8}
\newcommand{\depDayfacBaton}{10.8\%}
\newcommand{\depDayfacDetBaton}{14.2\%}
\newcommand{\depDayfacDetCf}{14.2\%}
\newcommand{\depDayfacDetComp}{two-lever}
\newcommand{\depDayfacDetCompSv}{11.8\%}
\newcommand{\depDayfacDetDelta}{+2.3}
\newcommand{\depDayfacDetDeltaCI}{[+1.4, +3.5]}
\newcommand{\depDayfacDetDpThree}{14.2\%}
\newcommand{\depDayfacDetHo}{9.8\%}
\newcommand{\depDayfacDetRatio}{100\%}
\newcommand{\depDayfacDetReact}{172.2}
\newcommand{\depDayfacDetThr}{9.1\%}
\newcommand{\depDayfacDetThrK}{9.7\%}
\newcommand{\depDayfacSaaBaton}{7.4\%}
\newcommand{\depDayfacSaaCf}{8.0\%}
\newcommand{\depDayfacSaaComp}{restock}
\newcommand{\depDayfacSaaCompSv}{5.8\%}
\newcommand{\depDayfacSaaDelta}{+1.6}
\newcommand{\depDayfacSaaDeltaCI}{[-1.6, +4.9]}
\newcommand{\depDayfacSaaDpThree}{12.4\%}
\newcommand{\depDayfacSaaHo}{-6.9\%}
\newcommand{\depDayfacSaaRatio}{60\%}
\newcommand{\depDayfacSaaReact}{5.2}
\newcommand{\depDayfacSaaThr}{-15.6\%}
\newcommand{\depDayfacSaaThrK}{-12.7\%}
\newcommand{\depNClose}{4}
\newcommand{\depRhoNineBaton}{49.9\%}
\newcommand{\depRhoNineDetBaton}{35.4\%}
\newcommand{\depRhoNineDetCf}{37.6\%}
\newcommand{\depRhoNineDetComp}{DP$^3_N$}
\newcommand{\depRhoNineDetCompSv}{36.5\%}
\newcommand{\depRhoNineDetDelta}{-1.1}
\newcommand{\depRhoNineDetDeltaCI}{[-1.6, -0.7]}
\newcommand{\depRhoNineDetDpThree}{38.0\%}
\newcommand{\depRhoNineDetHo}{34.1\%}
\newcommand{\depRhoNineDetRatio}{93\%}
\newcommand{\depRhoNineDetReact}{177.5}
\newcommand{\depRhoNineDetThr}{32.3\%}
\newcommand{\depRhoNineDetThrK}{34.0\%}
\newcommand{\depRhoNineSaaBaton}{64.5\%}
\newcommand{\depRhoNineSaaCf}{65.9\%}
\newcommand{\depRhoNineSaaComp}{two-lever}
\newcommand{\depRhoNineSaaCompSv}{54.4\%}
\newcommand{\depRhoNineSaaDelta}{+10.2}
\newcommand{\depRhoNineSaaDeltaCI}{[+8.8, +11.6]}
\newcommand{\depRhoNineSaaDpThree}{67.7\%}
\newcommand{\depRhoNineSaaHo}{33.0\%}
\newcommand{\depRhoNineSaaRatio}{95\%}
\newcommand{\depRhoNineSaaReact}{25.9}
\newcommand{\depRhoNineSaaThr}{27.3\%}
\newcommand{\depRhoNineSaaThrK}{32.9\%}
\newcommand{\depRhoSixBaton}{40.8\%}
\newcommand{\depRhoSixDetBaton}{27.9\%}
\newcommand{\depRhoSixDetCf}{29.2\%}
\newcommand{\depRhoSixDetComp}{DP$^3_N$}
\newcommand{\depRhoSixDetCompSv}{27.8\%}
\newcommand{\depRhoSixDetDelta}{+0.0}
\newcommand{\depRhoSixDetDeltaCI}{[-0.4, +0.4]}
\newcommand{\depRhoSixDetDpThree}{29.5\%}
\newcommand{\depRhoSixDetHo}{26.1\%}
\newcommand{\depRhoSixDetRatio}{94\%}
\newcommand{\depRhoSixDetReact}{179.2}
\newcommand{\depRhoSixDetThr}{24.7\%}
\newcommand{\depRhoSixDetThrK}{26.0\%}
\newcommand{\depRhoSixSaaBaton}{53.7\%}
\newcommand{\depRhoSixSaaCf}{53.8\%}
\newcommand{\depRhoSixSaaComp}{two-lever}
\newcommand{\depRhoSixSaaCompSv}{43.9\%}
\newcommand{\depRhoSixSaaDelta}{+9.7}
\newcommand{\depRhoSixSaaDeltaCI}{[+8.3, +11.3]}
\newcommand{\depRhoSixSaaDpThree}{56.4\%}
\newcommand{\depRhoSixSaaHo}{20.1\%}
\newcommand{\depRhoSixSaaRatio}{95\%}
\newcommand{\depRhoSixSaaReact}{16.5}
\newcommand{\depRhoSixSaaThr}{16.3\%}
\newcommand{\depRhoSixSaaThrK}{19.8\%}
\newcommand{\depRhoThreeBaton}{29.1\%}
\newcommand{\depRhoThreeDetBaton}{21.9\%}
\newcommand{\depRhoThreeDetCf}{22.5\%}
\newcommand{\depRhoThreeDetComp}{DP$^3_N$}
\newcommand{\depRhoThreeDetCompSv}{20.7\%}
\newcommand{\depRhoThreeDetDelta}{+1.2}
\newcommand{\depRhoThreeDetDeltaCI}{[+0.9, +1.6]}
\newcommand{\depRhoThreeDetDpThree}{22.7\%}
\newcommand{\depRhoThreeDetHo}{19.3\%}
\newcommand{\depRhoThreeDetRatio}{97\%}
\newcommand{\depRhoThreeDetReact}{180.2}
\newcommand{\depRhoThreeDetThr}{18.0\%}
\newcommand{\depRhoThreeDetThrK}{19.1\%}
\newcommand{\depRhoThreeSaaBaton}{36.3\%}
\newcommand{\depRhoThreeSaaCf}{36.6\%}
\newcommand{\depRhoThreeSaaComp}{restock}
\newcommand{\depRhoThreeSaaCompSv}{29.8\%}
\newcommand{\depRhoThreeSaaDelta}{+6.5}
\newcommand{\depRhoThreeSaaDeltaCI}{[+4.7, +8.4]}
\newcommand{\depRhoThreeSaaDpThree}{40.9\%}
\newcommand{\depRhoThreeSaaHo}{7.9\%}
\newcommand{\depRhoThreeSaaRatio}{89\%}
\newcommand{\depRhoThreeSaaReact}{8.0}
\newcommand{\depRhoThreeSaaThr}{3.1\%}
\newcommand{\depRhoThreeSaaThrK}{6.5\%}
\newcommand{\depRhoZeroBaton}{7.7\%}
\newcommand{\depRhoZeroDetBaton}{15.5\%}
\newcommand{\depRhoZeroDetCf}{15.4\%}
\newcommand{\depRhoZeroDetComp}{DP$^3_N$}
\newcommand{\depRhoZeroDetCompSv}{13.5\%}
\newcommand{\depRhoZeroDetDelta}{+2.0}
\newcommand{\depRhoZeroDetDeltaCI}{[+1.8, +2.1]}
\newcommand{\depRhoZeroDetDpThree}{15.5\%}
\newcommand{\depRhoZeroDetHo}{10.6\%}
\newcommand{\depRhoZeroDetRatio}{100\%}
\newcommand{\depRhoZeroDetReact}{176.4}
\newcommand{\depRhoZeroDetThr}{9.9\%}
\newcommand{\depRhoZeroDetThrK}{10.5\%}
\newcommand{\depRhoZeroSaaBaton}{-0.1\%}
\newcommand{\depRhoZeroSaaCf}{0.2\%}
\newcommand{\depRhoZeroSaaComp}{DP$_N$}
\newcommand{\depRhoZeroSaaCompSv}{0.0\%}
\newcommand{\depRhoZeroSaaDelta}{-0.1}
\newcommand{\depRhoZeroSaaDeltaCI}{[-4.5, +3.2]}
\newcommand{\depRhoZeroSaaDpThree}{7.4\%}
\newcommand{\depRhoZeroSaaHo}{-8.0\%}
\newcommand{\depRhoZeroSaaRatio}{-2\%}
\newcommand{\depRhoZeroSaaReact}{2.8}
\newcommand{\depRhoZeroSaaThr}{-24.9\%}
\newcommand{\depRhoZeroSaaThrK}{-11.9\%}
\newcommand{\detActGain}{1.8}
\newcommand{\dtAllAware}{27.6\%}
\newcommand{\dtAllPooled}{27.6\%}
\newcommand{\dtAwarePooledCI}{[-0.8, -0.3]}
\newcommand{\dtAwarePooledDelta}{-0.5}
\newcommand{\dtNNormal}{803}
\newcommand{\dtNPromo}{197}
\newcommand{\dtPromoAware}{29.7\%}
\newcommand{\dtPromoPooled}{29.4\%}
\newcommand{\dtPromoStale}{28.3\%}
\newcommand{\dtStalePooledCI}{[-1.7, -0.9]}
\newcommand{\dtStalePooledDelta}{-1.3}
\newcommand{\exDetBaton}{12.6\%}
\newcommand{\exDetBatonHo}{9.0\%}
\newcommand{\exDetDpThree}{12.6\%}
\newcommand{\exDetHo}{9.6\%}
\newcommand{\exDetMcErr}{0.39\%}
\newcommand{\exDetShare}{94\%}
\newcommand{\exDetShareDp}{94\%}
\newcommand{\exDetShareHo}{93\%}
\newcommand{\exDetThree}{13.4\%}
\newcommand{\exSaaBaton}{5.9\%}
\newcommand{\exSaaBatonHo}{-2.3\%}
\newcommand{\exSaaDpThree}{8.6\%}
\newcommand{\exSaaHo}{0.4\%}
\newcommand{\exSaaMcErr}{0.22\%}
\newcommand{\exSaaShare}{59\%}
\newcommand{\exSaaShareDp}{86\%}
\newcommand{\exSaaShareHo}{-607\%}
\newcommand{\exSaaThree}{10.0\%}
\newcommand{\feeBatonMax}{29.0\%}
\newcommand{\feeBreakEven}{29}
\newcommand{\feeGapMax}{0.0}
\newcommand{\feeGapZero}{19.1}
\newcommand{\feeHoMax}{29.0\%}
\newcommand{\feeMax}{30}
\newcommand{\frCityBiasIn}{-1.18\%}
\newcommand{\frCityBiasOut}{-1.05\%}
\newcommand{\frCityCf}{9.4\%}
\newcommand{\frCityDpThree}{10.7\%}
\newcommand{\frCityExact}{9.9\%}
\newcommand{\frCityFstar}{5.9\%}
\newcommand{\frCityHo}{9.2\%}
\newcommand{\frCityRsBaton}{0.8\%}
\newcommand{\frCityRsExact}{2.1\%}
\newcommand{\frCityThree}{5.6\%}
\newcommand{\frDethBiasIn}{+0.00\%}
\newcommand{\frDethBiasOut}{+0.01\%}
\newcommand{\frDethCf}{52.2\%}
\newcommand{\frDethDpThree}{54.6\%}
\newcommand{\frDethExact}{54.6\%}
\newcommand{\frDethFstar}{51.4\%}
\newcommand{\frDethHo}{18.9\%}
\newcommand{\frDethRsBaton}{4.1\%}
\newcommand{\frDethRsExact}{4.1\%}
\newcommand{\frDethThree}{51.4\%}
\newcommand{\frSalhiBiasIn}{+0.50\%}
\newcommand{\frSalhiBiasOut}{+0.09\%}
\newcommand{\frSalhiCf}{53.1\%}
\newcommand{\frSalhiDpThree}{53.5\%}
\newcommand{\frSalhiExact}{57.9\%}
\newcommand{\frSalhiFstar}{51.9\%}
\newcommand{\frSalhiHo}{41.3\%}
\newcommand{\frSalhiRsBaton}{20.1\%}
\newcommand{\frSalhiRsExact}{28.3\%}
\newcommand{\frSalhiThree}{52.2\%}
\newcommand{\lgCityBaton}{10.6\%}
\newcommand{\lgCityComp}{thr.-$k$}
\newcommand{\lgCityCompSv}{10.1\%}
\newcommand{\lgCityDelta}{+0.5}
\newcommand{\lgCityDeltaCI}{[+0.3, +0.7]}
\newcommand{\lgCityRatio}{86\%}
\newcommand{\lgCityRatioOrcThree}{28\%}
\newcommand{\lgCityUBaton}{11.8\%}
\newcommand{\lgCityUComp}{roll.-$\theta$}
\newcommand{\lgCityUCompSv}{11.7\%}
\newcommand{\lgCityUDelta}{+0.1}
\newcommand{\lgCityUDeltaCI}{[-0.7, +0.7]}
\newcommand{\lgCityURatio}{90\%}
\newcommand{\lgCityURatioOrcThree}{32\%}
\newcommand{\lgSalhiBaton}{53.0\%}
\newcommand{\lgSalhiComp}{DP$^3_N$}
\newcommand{\lgSalhiCompSv}{52.3\%}
\newcommand{\lgSalhiDelta}{+0.7}
\newcommand{\lgSalhiDeltaCI}{[+0.2, +1.2]}
\newcommand{\lgSalhiRatio}{98\%}
\newcommand{\lgSalhiRatioOrcThree}{85\%}
\newcommand{\lgZpFiftyBaton}{3.1\%}
\newcommand{\lgZpFiftyComp}{DP$_N$}
\newcommand{\lgZpFiftyCompSv}{3.1\%}
\newcommand{\lgZpFiftyDelta}{+0.0}
\newcommand{\lgZpFiftyDeltaCI}{[-6.9, +5.0]}
\newcommand{\lgZpFiftyRatio}{47\%}
\newcommand{\lgZpFiftyRatioOrcThree}{8\%}
\newcommand{\lgZpTwentyFiveBaton}{4.9\%}
\newcommand{\lgZpTwentyFiveComp}{thr.-$k$}
\newcommand{\lgZpTwentyFiveCompSv}{4.5\%}
\newcommand{\lgZpTwentyFiveDelta}{+0.4}
\newcommand{\lgZpTwentyFiveDeltaCI}{[+0.2, +0.6]}
\newcommand{\lgZpTwentyFiveRatio}{76\%}
\newcommand{\lgZpTwentyFiveRatioOrcThree}{15\%}
\newcommand{\nBatonTop}{six}
\newcommand{\nCfAbove}{all}
\newcommand{\nCity}{19}
\newcommand{\nOrcBeat}{four}
\newcommand{\nRL}{49}
\newcommand{\nRegretRoutes}{735}
\newcommand{\nTimingRoutes}{102}
\newcommand{\poolBatonMax}{7}
\newcommand{\poolBatonMed}{0}
\newcommand{\poolCityBeatLow}{0/19}
\newcommand{\poolCityBreakEven}{2.6}
\newcommand{\poolCityGainOne}{5.4\%}
\newcommand{\poolCityGainZero}{10.7\%}
\newcommand{\poolCityK}{8}
\newcommand{\poolCityNaiveOne}{36.0}
\newcommand{\poolCityNaiveZero}{41.1}
\newcommand{\poolCityPpu}{33.4}
\newcommand{\poolCityShadowOne}{34.1}
\newcommand{\poolCityShadowZero}{36.7}
\newcommand{\poolCitySstar}{0.0}
\newcommand{\poolCitySstarLow}{0.1}
\newcommand{\poolCityTotHigh}{36.7}
\newcommand{\poolCityTotLow}{36.6}
\newcommand{\poolCuiMax}{1}
\newcommand{\poolCuiMed}{0}
\newcommand{\poolDetMax}{7}
\newcommand{\poolDetMed}{3.5}
\newcommand{\poolDethBeatLow}{40/40}
\newcommand{\poolDethBreakEven}{19.9}
\newcommand{\poolDethGainOne}{4.2\%}
\newcommand{\poolDethGainZero}{5.0\%}
\newcommand{\poolDethK}{7}
\newcommand{\poolDethNaiveOne}{154.3}
\newcommand{\poolDethNaiveZero}{176.5}
\newcommand{\poolDethPpu}{135.0}
\newcommand{\poolDethSaaBeatLow}{21/40}
\newcommand{\poolDethSaaBreakEven}{0.1}
\newcommand{\poolDethSaaGainOne}{0.6\%}
\newcommand{\poolDethSaaGainZero}{0.3\%}
\newcommand{\poolDethSaaK}{10}
\newcommand{\poolDethSaaNaiveOne}{8.2}
\newcommand{\poolDethSaaNaiveZero}{8.3}
\newcommand{\poolDethSaaPpu}{8.3}
\newcommand{\poolDethSaaShadowOne}{8.1}
\newcommand{\poolDethSaaShadowZero}{8.3}
\newcommand{\poolDethSaaSstar}{0.0}
\newcommand{\poolDethSaaSstarLow}{0.0}
\newcommand{\poolDethSaaTotHigh}{8.3}
\newcommand{\poolDethSaaTotLow}{8.3}
\newcommand{\poolDethShadowOne}{147.8}
\newcommand{\poolDethShadowZero}{167.7}
\newcommand{\poolDethSstar}{1.1}
\newcommand{\poolDethSstarLow}{4.2}
\newcommand{\poolDethTotHigh}{162.5}
\newcommand{\poolDethTotLow}{119.0}
\newcommand{\poolGounarisMax}{1}
\newcommand{\poolGounarisMed}{0}
\newcommand{\poolHoMax}{8}
\newcommand{\poolHoMed}{1}
\newcommand{\poolMDROMax}{0}
\newcommand{\poolMDROMed}{0}
\newcommand{\poolMetroBeatLow}{4/4}
\newcommand{\poolMetroBreakEven}{6.6}
\newcommand{\poolMetroGainMax}{9.4\%}
\newcommand{\poolMetroGainMaxS}{19}
\newcommand{\poolMetroGainOne}{3.0\%}
\newcommand{\poolMetroGainZero}{5.0\%}
\newcommand{\poolMetroK}{66}
\newcommand{\poolMetroLamLarge}{3.3}
\newcommand{\poolMetroLamSmall}{40}
\newcommand{\poolMetroNaiveOne}{1721.7}
\newcommand{\poolMetroNaiveZero}{1765.4}
\newcommand{\poolMetroPpu}{1349.9}
\newcommand{\poolMetroShadowOne}{1670.7}
\newcommand{\poolMetroShadowZero}{1677.3}
\newcommand{\poolMetroSmax}{40}
\newcommand{\poolMetroSstar}{15.5}
\newcommand{\poolMetroSstarLow}{27.2}
\newcommand{\poolMetroTotHigh}{1421.5}
\newcommand{\poolMetroTotLow}{1109.2}
\newcommand{\poolSAAMax}{0}
\newcommand{\poolSAAMed}{0}
\newcommand{\poolWDROMax}{0}
\newcommand{\poolWDROMed}{0}
\newcommand{\ratioAllHi}{100\%}
\newcommand{\ratioAllLo}{47\%}
\newcommand{\ratioDpThreeHi}{96\%}
\newcommand{\ratioDpThreeLo}{86\%}
\newcommand{\ratioOrcThreeHi}{68\%}
\newcommand{\ratioOrcThreeLo}{52\%}
\newcommand{\regretAll}{3.8\%}
\newcommand{\regretBoundAll}{7.7\%}
\newcommand{\regretBoundFail}{1}
\newcommand{\regretBoundHolds}{734/735}
\newcommand{\regretCity}{2.0\%}
\newcommand{\regretCleanShare}{90\%}
\newcommand{\regretContained}{100.0\%}
\newcommand{\regretDet}{3.8\%}
\newcommand{\regretFailMax}{0.06\%}
\newcommand{\regretLong}{4.5\%}
\newcommand{\regretMid}{6.8\%}
\newcommand{\regretMin}{2.0\%}
\newcommand{\regretMinRow}{City \textsc{Det}}
\newcommand{\regretRouteMean}{5.1\%}
\newcommand{\regretSaa}{6.8\%}
\newcommand{\regretShort}{2.7\%}
\newcommand{\rlBaton}{33.7\%}
\newcommand{\rlBatonFitS}{1.5}
\newcommand{\rlBatonHo}{32.9\%}
\newcommand{\rlBest}{27.4\%}
\newcommand{\rlInstCI}{[4.4, 7.7]}
\newcommand{\rlInstDelta}{6.0}
\newcommand{\rlInstWin}{12/12}
\newcommand{\rlNInst}{12}
\newcommand{\rlTrainMax}{3}
\newcommand{\rlTrainMin}{1}
\newcommand{\rlWinN}{39/49}
\newcommand{\rlWinP}{$p \le 10^{-3}$}
\newcommand{\salhiActShare}{89\%}
\newcommand{\sbSixtyBaton}{36.8\%}
\newcommand{\sbSixtyHo}{2.1\%}
\newcommand{\sbSixtyThresh}{1.4\%}
\newcommand{\shapePowTenAll}{64\%}
\newcommand{\shapePowTenDayfac}{47\%}
\newcommand{\shapePowTenNine}{88\%}
\newcommand{\shapePowTenSix}{72\%}
\newcommand{\shapePowTenThree}{58\%}
\newcommand{\shapePowTenZero}{51\%}
\newcommand{\shapePowTwentyFiveAll}{89\%}
\newcommand{\shapePowTwentyFiveDayfac}{78\%}
\newcommand{\shapePowTwentyFiveNine}{97\%}
\newcommand{\shapePowTwentyFiveSix}{97\%}
\newcommand{\shapePowTwentyFiveThree}{89\%}
\newcommand{\shapePowTwentyFiveZero}{82\%}
\newcommand{\shapeViolDayfac}{0.6\%}
\newcommand{\shapeViolNine}{0.0\%}
\newcommand{\shapeViolSix}{0.1\%}
\newcommand{\shapeViolThree}{0.1\%}
\newcommand{\shapeViolZero}{0.2\%}
\newcommand{\shopTravelSaving}{25.9\%}
\newcommand{\statCompDeltaMax}{10.3}
\newcommand{\statCompDeltaMin}{0.0}
\newcommand{\statCompLoMin}{-0.4}
\newcommand{\statCompWinMin}{21}
\newcommand{\statConsDeltaMax}{10.3}
\newcommand{\statConsDeltaMin}{6.5}
\newcommand{\statConsPmax}{$p_{\mathrm{Holm}} < 10^{-3}$}
\newcommand{\statConsWinMin}{31}
\newcommand{\statCuiCI}{[+8.1, +12.6]}
\newcommand{\statCuiComp}{two-lever}
\newcommand{\statCuiDelta}{+10.3}
\newcommand{\statCuiDpDelta}{-2.5}
\newcommand{\statCuiP}{0.00}
\newcommand{\statCuiWin}{40}
\newcommand{\statDetCI}{[-0.4, +0.4]}
\newcommand{\statDetComp}{DP$^3_N$}
\newcommand{\statDetDelta}{+0.0}
\newcommand{\statDetDpDelta}{-1.6}
\newcommand{\statDetP}{0.49}
\newcommand{\statDetWin}{21}
\newcommand{\statDpDeltaMax}{-1.6}
\newcommand{\statDpDeltaMin}{-6.6}
\newcommand{\statGounarisCI}{[+7.7, +11.2]}
\newcommand{\statGounarisComp}{two-lever}
\newcommand{\statGounarisDelta}{+9.4}
\newcommand{\statGounarisDpDelta}{-3.1}
\newcommand{\statGounarisP}{0.00}
\newcommand{\statGounarisWin}{40}
\newcommand{\statMDROCI}{[+5.1, +11.2]}
\newcommand{\statMDROComp}{restock}
\newcommand{\statMDRODelta}{+7.9}
\newcommand{\statMDRODpDelta}{-5.0}
\newcommand{\statMDROP}{0.00}
\newcommand{\statMDROWin}{34}
\newcommand{\statPmax}{$p_{\mathrm{Holm}} \le 0.489$}
\newcommand{\statSAACI}{[+8.3, +11.2]}
\newcommand{\statSAAComp}{two-lever}
\newcommand{\statSAADelta}{+9.7}
\newcommand{\statSAADpDelta}{-2.7}
\newcommand{\statSAAP}{0.00}
\newcommand{\statSAAWin}{40}
\newcommand{\statWDROCI}{[+3.6, +9.7]}
\newcommand{\statWDROComp}{restock}
\newcommand{\statWDRODelta}{+6.5}
\newcommand{\statWDRODpDelta}{-6.6}
\newcommand{\statWDROP}{0.00}
\newcommand{\statWDROWin}{31}
\newcommand{\svCityBaton}{10.6\%}
\newcommand{\svCityCf}{10.8\%}
\newcommand{\svCityDp}{12.1\%}
\newcommand{\svCityHo}{10.6\%}
\newcommand{\svCityOrc}{36.8\%}
\newcommand{\svCuiBaton}{52.5\%}
\newcommand{\svCuiCf}{53.0\%}
\newcommand{\svCuiDpN}{13.0\%}
\newcommand{\svCuiDpThree}{54.9\%}
\newcommand{\svCuiDpThreeN}{34.3\%}
\newcommand{\svCuiHo}{20.7\%}
\newcommand{\svCuiOrc}{43.1\%}
\newcommand{\svCuiOrcThree}{76.6\%}
\newcommand{\svCuiRoTheta}{18.7\%}
\newcommand{\svCuiThrK}{20.3\%}
\newcommand{\svCuiThresh}{16.7\%}
\newcommand{\svCuiTwo}{42.2\%}
\newcommand{\svDetBaton}{27.9\%}
\newcommand{\svDetCf}{29.2\%}
\newcommand{\svDetDpN}{25.6\%}
\newcommand{\svDetDpThree}{29.5\%}
\newcommand{\svDetDpThreeN}{27.8\%}
\newcommand{\svDetHo}{26.1\%}
\newcommand{\svDetOrc}{40.7\%}
\newcommand{\svDetOrcThree}{47.2\%}
\newcommand{\svDetRoTheta}{24.9\%}
\newcommand{\svDetThrK}{26.0\%}
\newcommand{\svDetThresh}{24.7\%}
\newcommand{\svDetTwo}{26.5\%}
\newcommand{\svGounarisBaton}{53.2\%}
\newcommand{\svGounarisCf}{53.8\%}
\newcommand{\svGounarisDpN}{14.6\%}
\newcommand{\svGounarisDpThree}{56.3\%}
\newcommand{\svGounarisDpThreeN}{36.2\%}
\newcommand{\svGounarisHo}{21.0\%}
\newcommand{\svGounarisOrc}{43.3\%}
\newcommand{\svGounarisOrcThree}{78.0\%}
\newcommand{\svGounarisRoTheta}{19.1\%}
\newcommand{\svGounarisThrK}{20.7\%}
\newcommand{\svGounarisThresh}{16.7\%}
\newcommand{\svGounarisTwo}{43.8\%}
\newcommand{\svMDROBaton}{45.3\%}
\newcommand{\svMDROCf}{45.6\%}
\newcommand{\svMDRODpN}{0.0\%}
\newcommand{\svMDRODpThree}{50.2\%}
\newcommand{\svMDRODpThreeN}{-4.7\%}
\newcommand{\svMDROHo}{11.8\%}
\newcommand{\svMDROOrc}{42.2\%}
\newcommand{\svMDROOrcThree}{79.4\%}
\newcommand{\svMDRORoTheta}{10.3\%}
\newcommand{\svMDROThrK}{8.5\%}
\newcommand{\svMDROThresh}{0.5\%}
\newcommand{\svMDROTwo}{35.6\%}
\newcommand{\svSAABaton}{53.7\%}
\newcommand{\svSAACf}{53.8\%}
\newcommand{\svSAADpN}{12.4\%}
\newcommand{\svSAADpThree}{56.4\%}
\newcommand{\svSAADpThreeN}{35.4\%}
\newcommand{\svSAAHo}{20.1\%}
\newcommand{\svSAAOrc}{43.5\%}
\newcommand{\svSAAOrcThree}{78.4\%}
\newcommand{\svSAARoTheta}{18.4\%}
\newcommand{\svSAAThrK}{19.8\%}
\newcommand{\svSAAThresh}{16.3\%}
\newcommand{\svSAATwo}{43.9\%}
\newcommand{\svSalhiBaton}{53.0\%}
\newcommand{\svSalhiDpThree}{54.0\%}
\newcommand{\svSalhiHo}{42.3\%}
\newcommand{\svSalhiOrc}{48.8\%}
\newcommand{\svWDROBaton}{41.2\%}
\newcommand{\svWDROCf}{41.7\%}
\newcommand{\svWDRODpN}{0.0\%}
\newcommand{\svWDRODpThree}{47.8\%}
\newcommand{\svWDRODpThreeN}{-1.4\%}
\newcommand{\svWDROHo}{10.0\%}
\newcommand{\svWDROOrc}{42.0\%}
\newcommand{\svWDROOrcThree}{79.4\%}
\newcommand{\svWDRORoTheta}{8.9\%}
\newcommand{\svWDROThrK}{7.2\%}
\newcommand{\svWDROThresh}{-1.3\%}
\newcommand{\svWDROTwo}{32.4\%}
\newcommand{\svZpFiftyBaton}{3.1\%}
\newcommand{\svZpFiftyThr}{-1.5\%}
\newcommand{\svZpTwentyFiveBaton}{4.9\%}
\newcommand{\tailBatonHi}{54.3\%}
\newcommand{\tailBatonLo}{22.3\%}
\newcommand{\tailDetBaton}{22.3\%}
\newcommand{\tailDetHo}{23.8\%}
\newcommand{\tailDetThr}{23.6\%}
\newcommand{\thrGapCity}{0.4\%}
\newcommand{\thrGapDet}{1.0\%}
\newcommand{\thrGapLong}{0.9\%}
\newcommand{\thrGapMid}{1.7\%}
\newcommand{\thrGapSaa}{3.0\%}
\newcommand{\thrGapShort}{0.8\%}
\newcommand{\thrkHoGapMax}{3.3}
\newcommand{\thrkHoGapMin}{0.1}
\newcommand{\timeBatonHoMs}{7\,ms}
\newcommand{\timeBatonMs}{24\,ms}
\newcommand{\timeOnlineSkUs}{63\,$\mu$s}
\newcommand{\timeOnlineUs}{1.2\,$\mu$s}
\newcommand{\twoGapHi}{10.3}
\newcommand{\twoGapLo}{1.4}


\setlength{\emergencystretch}{2em}
\definecolor{rev}{RGB}{70,70,70}
\newenvironment{comment}{\par\medskip\noindent\color{rev}\itshape
  \begin{list}{}{\leftmargin=1.2em\rightmargin=1.2em}\item[]}%
  {\end{list}\par}
\newcommand{\resp}{\par\noindent\textbf{Response.}\ }
\newcommand{\chg}{\par\noindent\textbf{Changes.}\ }
\newcommand{\sref}[1]{Section~\ref{m-#1}}
\newcommand{\tref}[1]{Table~\ref{m-#1}}
\newcommand{\fref}[1]{Figure~\ref{m-#1}}
\newcommand{\pref}[1]{Proposition~\ref{m-#1}}
\newcommand{\BATON}{\textsc{Baton}}

\title{Response to the reviewers\\[4pt]
\large Manuscript CAOR-D-26-01885, ``Optimal-stopping recourse for
vehicle routing with simultaneous pickup, delivery, and stochastic
demands''}
\author{Quang-Vinh Dang, Hoang-Viet Vu, Phuc-Son Nguyen}
\date{}

\begin{document}
\maketitle

\noindent Dear Editor, dear Reviewers,

We thank the three reviewers for careful and constructive reports. They
asked us to clarify the stochastic assumptions and the depot-return
model, to test the method under dependent and non-exchangeable demand
and under a shared standby pool, to compare against stronger threshold
policies, to report comparable data and computing budgets, to
strengthen the statistical analysis, and to state the contribution and
the scope of the conclusions more carefully. We have addressed each of
these points. Several of the new experiments did not come out entirely
in our favour, and we report those results in the manuscript as they
are: an equal-data three-action dynamic program ties \BATON{} on the
deterministic plans and beats it under very strong demand correlation,
and on two city variants \BATON{} is only tied with its
strongest competitor. Every comment is answered below; the reviewers' text is set in italics,
and section, table and proposition numbers refer to the revised
manuscript. A version of the manuscript with all changes marked against
the submission is provided separately.

\section*{Summary of the main changes}

\begin{enumerate}[leftmargin=1.4em]
\item \textbf{Assumptions and theory} (\sref{sec:formulation},
  \sref{sec:defects}, \sref{sec:estimation}). Assumption~1 now requires
  only that the load be a stochastically monotone Markov chain, which
  covers independent increments and a Gaussian day-factor model;
  Assumption~2 now requires only that the emergency price be
  non-increasing along the route---no ordering between handoff and
  emergency prices is needed. Proposition~\ref{m-prop:monotone}
  (monotone continuation value) is proved under these weaker assumptions
  and for the three-action recursion. A new
  Proposition~\ref{m-prop:regret} quantifies the cost of over-triggering
  (Reviewer~2, comment~4).
\item \textbf{New experiments} (\sref{sec:robust}--\sref{sec:budget}):
  demand correlation $\rho \in \{0, 0.3, 0.6, 0.9\}$ and a day-level
  demand factor, with a direct test of monotonicity and of its power;
  an exact dynamic program under independent demand; promotion days and
  day-type-specific fits; the bias of the fresh-start value and the
  exact post-return state; a capacitated standby pool shared by the
  routes of a plan or of a metropolitan group of plans, with a sweep of
  holding costs and a shadow price;
  training-data budgets from 100 to 20{,}000 days; single-core timing of
  every policy; deliver-only customers; city plans re-planned under the
  \textsc{SAA} gate; a configuration in which standby is dearer than an
  emergency; and a fixed fee on every depot return, with the fee at which
  the return lever stops paying (\feeBreakEven{} currency units).
\item \textbf{New competitors and references}: a position-dependent
  threshold (one tuned level per stop), a cost-scaled rollout (a
  cost-dependent threshold with a tuned price multiplier), a
  \emph{two-lever} rule that gives a tuned threshold the depot return as
  well, and an equal-data dynamic program for the \emph{three-action}
  problem (DP$^3_N$), so that \BATON{} is compared with policies that
  have its own action set. As references we added a high-data plug-in
  program for the three-action problem, the three-action clairvoyant
  bound, and an \emph{exact} dynamic program, solved by numerical
  convolution, for the case of independent demand
  (\tref{tab:exact}).
\item \textbf{Statistics}: the instance within a gate is now the
  experimental unit (the routes of a plan share their days), with paired
  mean differences, 95\% bootstrap confidence intervals, win counts,
  effect sizes and Holm-adjusted tests per gate (\tref{tab:stats}),
  paired intervals for the key contrast of every robustness table,
  tail-risk measures (\tref{tab:tail}) and route-level results as
  supplementary data.
\item \textbf{Positioning and scope}: the introduction separates the
  modelling contribution, the structural results and the adaptation of
  established regression-based stopping methods; the conclusion states
  the scope of the managerial conclusions and the boundaries of the
  action set.
\item \textbf{Length}: five supplementary tables, one figure and the
  proof of Proposition~\ref{m-prop:bias} have moved to an appendix.
\item \textbf{Machine-checked proofs}: all four propositions have been
  verified formally in the Lean~4 proof assistant with the Mathlib
  library. The formalisation compiles without errors or unproved steps
  and relies only on Lean's standard axioms. It is released with the
  replication code (directory \texttt{papers/baton/BatonProofs}), and its
  correspondence to the text is documented there. The formalisation
  abstracts the transition law to a kernel with the monotonicity
  property, and proves Propositions~\ref{m-prop:myopic}
  and~\ref{m-prop:regret} on finite trees of histories; the manuscript
  now lists these abstractions (\sref{sec:method}, first paragraph).
\end{enumerate}

In the course of the revision we also found and corrected several
errors of our own; they are listed at the end of this letter. The most
important is that the proof of the cost bound in the original
Proposition~1 compared the optimal policy with the clairvoyant bound in
the wrong direction. The corrected Proposition~\ref{m-prop:bias} gives
two-sided bounds that are now proved.

% ════════════════════════════════════════════════════════════════════════
\section*{Reviewer 1}

\begin{comment}
Stochastic assumptions and state sufficiency. The manuscript should
clarify the independence assumptions required for the theoretical
results. In particular, the authors should justify under which
conditions the route position and accumulated load constitute a
sufficient state. An experiment with correlated demands or route-level
demand factors would strengthen the analysis.
\end{comment}
\resp We agree that the original Assumption~1 (independent increments)
was stronger than needed, and that its relation to the correlated
experiments was not explained. The revised Assumption~1 asks that the
running load $(W_k)$ be a Markov chain whose transition kernels are
stochastically monotone. Under this condition $(k, W_k)$ is a sufficient
state for the handoff problem and all results go through; for the
three-action problem it is sufficient under independent increments, and
the manuscript now says after Eq.~\eqref{m-eq:bellman3} that the
recursion is misspecified under factor dependence, which the
state-conditional variant \textsc{Baton-cf} corrects. The new text identifies two cases of
practical interest: (i) independent, not necessarily identically
distributed increments, and (ii) positively dependent demand driven by
a common day-level factor, $g_j = \mu_j + \beta Z + \varepsilon_j$ with
Gaussian $Z$ and noise of common variance. In case (ii) the likelihood
depends on $Z$ only through $W_k$, so $W_k$ is sufficient for the
factor, its posterior is increasing in $W_k$, and the chain is Markov
with monotone kernels. Outside these cases---including our benchmark
protocol, where pickups and deliveries load on separate factors with
customer-specific scales---$W_k$ is not exactly sufficient, and we now
say so and describe the policy as the regression-based approximation
within the class of policies that observe $(k,W_k)$.

We then measured what this approximation costs. \tref{tab:dependence}
re-fits and re-tests every policy on the Dethloff \textsc{Det} and
\textsc{SAA} plans under copula correlations $\rho = 0, 0.3, 0.6, 0.9$
and under a route-level multiplicative day factor. Dependence turns out
to help rather than hurt: every implementable policy saves more as $\rho$ grows,
because correlated demand creates heavy days whose breaches can be
foreseen from the running load. A \BATON{} variant has the highest
saving among the implementable policies under every law. \BATON{}
itself leads the strongest competitor with a paired interval above zero
in six of the ten rows and is tied in three; in one row it loses: at
$\rho = 0.9$ on \textsc{Det} plans the equal-data three-action dynamic
program is ahead by \depRhoNineDetDelta{} points
(\depRhoNineDetDeltaCI). That is where the dependence bias of the
fresh-start value is strongest, and \textsc{Baton-cf} reverses it
(\depRhoNineDetCf{} against \depRhoNineDetCompSv). On the \textsc{Det}
plans \BATON{} attains \depRhoZeroDetRatio, \depRhoSixDetRatio{} and
\depRhoNineDetRatio{} of the high-data three-action saving at
$\rho = 0$, $0.6$ and $0.9$. Under independence we can also compute the
exact optimum by numerical convolution (\tref{tab:exact}): on
\textsc{Det} plans \BATON{} attains \exDetShare{} of the exact
three-action saving. We also tested the shape directly: the
unconstrained binned estimate of $\mathbb E[\text{future cost}\mid W_k]$
on fifty thousand paths per route decreases significantly between
adjacent bins in \shapeViolZero{} of pairs under independence, where
monotonicity is a theorem and this rate measures the false-positive
level of the test, in \shapeViolThree, \shapeViolSix{} and
\shapeViolNine{} of pairs at $\rho = 0.3$, $0.6$ and $0.9$, and in
\shapeViolDayfac{} under the day factor. The test found no significant
violations beyond its false-positive level, but its power is limited: a
dip of 10\% (25\%) of the mean continuation cost injected into one bin
is detected in only \shapePowTenAll{} (\shapePowTwentyFiveAll) of the
tests. The evidence is therefore consistent with a monotone continuation
value, without proving it, and the manuscript says so.

The main place where dependence matters is a finding we did not
anticipate. It concerns the depot return: under dependence the value of
the post-return state depends on the state that triggered the return,
and an unconditional fresh-start value misprices it on exactly the
days on which it is used. This is discussed under Reviewer~2,
comment~2, and a state-conditional variant (\textsc{Baton-cf}) removes
the effect.
\chg Revised Assumption~1 and the paragraph after it
(\sref{sec:formulation}); proof of Proposition~\ref{m-prop:monotone}
(\sref{sec:estimation}); new \sref{sec:robust} with
\tref{tab:dependence}.

\begin{comment}
Shared standby capacity. The policy appears to be applied independently
to each route, although standby vehicles are described as a pooled
resource. If standby capacity is limited, decisions across routes are
coupled. The authors should clarify whether capacity is unlimited,
allocated in advance, or represented through an exogenous shadow price.
\end{comment}
\resp In the original model standby capacity is pay-per-use: the pool
is shared across a metropolitan operation and large relative to one
plan, and each handoff is billed one standby vehicle-day
($F_{\mathrm{sb}}$), so no plan exhausts it. We now state this
explicitly and report the pool size each policy would need (the 95th
percentile of daily handoffs per plan: for \BATON{} a median of
\poolDetMed{} and a maximum of \poolDetMax{} vehicles on the Dethloff
\textsc{Det} plans, and a median of \poolSAAMed{} on the \textsc{SAA}
plans).

We agree that the capacitated case is important and have added it as a
new experiment (\sref{sec:pool}, \tref{tab:pool}). A pool of $S$
standby vehicles is reserved for the day, so a handoff then costs only
its marginal price, and the pool serves either the routes of one plan
or, in a metropolitan setting, the \textsc{Det} plans of all ten
instances of a Dethloff class (about \poolMetroK{} routes). Requests are
served first come, first served in the order in which they arise along
the routes, and a route whose request is refused decides at the refusal
stop between continuing and a depot return. We compare a naive policy
that ignores the cap with a policy that decides each route with a shadow
price $\lambda$ on the pool, chosen on the training days, and with the
two fixed prices $\lambda = F_{\mathrm{sb}}$ and $\lambda = \infty$; we
report the break-even holding cost of a reserved vehicle and the
cost-minimising pool size at holding costs of 20 and 5 per vehicle-day.
Three results follow. First, a reserved vehicle serves at most one route
a day, so at a holding cost equal to the standby day rate a reserved pool
cannot beat pay-per-use, and in our data it never does; below that rate,
reservation pays where handoffs are frequent. At a holding cost of 5 the
optimal reserved pool beats pay-per-use on \poolDethBeatLow{} Dethloff
\textsc{Det} plans (\poolDethTotLow{} against \poolDethPpu{} per plan-day)
and on \poolMetroBeatLow{} metropolitan pools, but not on the slack
\textsc{SAA} plans or the city plans, which need few handoffs. Second,
given a cap the routes are coupled, and deciding each route with a
shadow price beats ignoring the cap wherever the cap binds (by
\poolDethGainOne{} on Dethloff \textsc{Det} plans with one reserved
vehicle, with a paired interval above zero, and by up to
\poolMetroGainMax{} on the metropolitan pools). Third, the shadow price
chosen on the training days is high for small pools and falls as the
pool grows, and neither fixed price matches it. The shadow price keeps
the routes decoupled, so the policy is still fitted one route at a time,
and the depot return is the lever that makes a small or empty pool
tolerable. The three options the reviewer lists are therefore all
covered: unlimited pooled capacity (the default model), capacity
allocated in advance (the reserved pool, with its size and break-even
cost), and an exogenous shadow price (the $\lambda$ that decouples the
routes).
\chg New paragraph ``Standby versus emergency vehicles'' in
\sref{sec:setting}; new \sref{sec:pool} with
\tref{tab:pool}; the conclusion lists shared resources among the scope
limits.

\begin{comment}
Depot-return transition. The state transition after returning to the
depot should be described more precisely. The manuscript should explain
the vehicle load after unloading and reloading, the resulting residual
capacity, whether multiple returns are allowed, and how this action
affects the continuation-value recursion.
\end{comment}
\resp Thank you; the original description was too brief and, in one
respect, imprecise. The revised \sref{sec:actions} now opens with a
paragraph on the transition. After stop $k$ the vehicle carries
$L_0 + W_k$. A return unloads all pickups collected so far and leaves
on board exactly the undelivered cargo $L_0 - D_{\le k}$. The policy
records the post-return state as $W = 0$ against the unchanged slack
$B$, which is a conservative convention: the exact post-return state is
$x_k = -D_{\le k}$, so the convention understates the residual capacity
by the volume already delivered. We adopt it because it fixes the
post-return state and keeps the fresh-start value a single number per
stop. Returns may be repeated: after a return the vehicle re-enters the
recursion in state $(k,0)$, so further returns and a later handoff
remain available, and the recursion acquires the single extra term
$R_k + C_k(0)$ in~\eqref{m-eq:bellman3}. To quantify the convention, we
implemented the exact reset, with a fresh-start value $F_k(x)$ that
depends on the observed delivered volume (\tref{tab:fresh}, column
``exact reset''). Without deployment selection, pricing the exact
residual capacity raises the three-action saving from \frSalhiThree{} to
\frSalhiExact{} on Salhi--Nagy and from \frDethThree{} to \frDethExact{}
on Dethloff \textsc{SAA} plans. The convention used in the main tables
therefore understates the value of the depot return, and the reported
results for the three-action policy are conservative. We kept the
conservative convention in the main tables to preserve comparability
with the submitted version and with the published restocking competitor,
which uses the same reset. We state the exact reset as the recommended
specification for deployment.
\chg New paragraph ``The depot-return transition'' in
\sref{sec:actions}; exact-reset results in \sref{sec:robust} and
\tref{tab:fresh}.

\begin{comment}
Threshold-policy comparison. The theoretical criticism applies mainly
to fixed global thresholds. More flexible policies, such as
position-dependent or cost-dependent thresholds, may remain
competitive. Including such a benchmark would help isolate the value of
estimating the full continuation function.
\end{comment}
\resp This is a good point, and we have added both, together with a
third benchmark that the question suggests. The
\emph{position-dependent threshold} (thr.-$k$) has one load threshold
per stop, tuned by coordinate descent on realised training cost.
Because $C_k$ is monotone, the optimal stopping region at each stop is
an upper set $\{w > w_k^\star\}$, so this class can represent the
optimal handoff boundary exactly, and it isolates precisely what the
reviewer asks: the value of \emph{estimating} the boundary by backward
induction rather than searching for it. The coordinate descent starts
from the best of the reactive policy, the per-stop cuts of the tuned
global threshold and the per-stop boundaries of \BATON-ho, so it is at
least as good as each of them on the training days. The
\emph{cost-scaled rollout} (roll.-$\theta$) hands off when the myopic
cost $C^0_k(W_k)$ exceeds $\theta H_k$ with a tuned $\theta$; this is a
cost-dependent threshold that uses the per-stop prices. Finally, the
\emph{two-lever} rule gives a tuned threshold the depot return as well,
with a return trigger that may fire repeatedly and both parameters tuned
jointly; together with the equal-data three-action dynamic program
DP$^3_N$, it separates the value of \emph{pricing} two levers from the
value of \emph{having} them.

Results (\tref{tab:grand}, \tref{tab:large}). For the handoff decision
alone the reviewer's expectation is borne out: the position-dependent
threshold comes within \thrkHoGapMin--\thrkHoGapMax{} points of
\BATON-ho on every gate (\svDetThrK{} against \svDetHo{} on
\textsc{Det}), so a well-started search recovers most of the handoff
boundary. The cost-scaled rollout behaves similarly. The two-lever rule
captures part of the depot-return value but stays
\twoGapLo--\twoGapHi{} points below \BATON{} across the gates, and
DP$^3_N$ ties \BATON{} on the deterministic plans but collapses on the
conservative ones, where breaches are rare and bins cannot pool
information along the load axis. On the five conservative gates
\BATON's advantage over the strongest of all these competitors is
\statConsDeltaMin--\statConsDeltaMax{} points with intervals above zero
(\tref{tab:stats}). We now also say explicitly, after
\pref{prop:myopic}, that the critique applies to the myopic break-even
rule, and that a stop-dependent threshold can represent the optimal
boundary but must then be searched.
\chg Three new competitors and DP$^3_N$ in \sref{sec:setup} and
\tref{tab:competitors}; new columns in \tref{tab:grand},
\tref{tab:large} and \tref{tab:dependence}; discussion in the first and
third observations of \sref{sec:headline} and after
\pref{prop:myopic}.

\begin{comment}
Computational fairness. The methods use different amounts of training
data and computational resources. The authors should report performance
under comparable data and computational budgets, clearly separating
offline training time from online decision time. The training,
validation, and testing protocol should also be fully documented.
\end{comment}
\resp We have (i) documented the protocol, (ii) made the data budgets
explicit and varied them, and (iii) measured offline and online time
separately. (i) The new ``Protocol'' paragraph in \sref{sec:setup}
states that every data-driven policy---\BATON, the thresholds, the
two-lever rule, the rule-based triggers, the rollouts, the restocking
trigger, DP$_N$, DP$^3_N$ and the RL policy---sees exactly the same
$10^3$ training days per instance.
Hyper-parameters of tuned competitors are selected on realised cost
over those same days; there is no separate validation set, which if
anything favours the tuned competitors, since \BATON{} has no
continuous hyper-parameter (its one discrete choice, the recourse menu,
is made on the same days). The $2\times10^3$ test days are used only for evaluation, and the
seeds are fixed functions of the instance name. The only policies with
more data are the two high-data reference programs, which are
yardsticks rather than competitors; every table that reports them labels
them, and the clairvoyant bounds, as reference points, and the main
tables separate them from the competitors. (ii) The new data-budget experiment (\sref{sec:budget},
\tref{tab:budget} and \fref{fig:budget}) varies the history from 100 to
20{,}000 days per route for \BATON, the tuned threshold and the
equal-data dynamic program. \BATON{} beats the tuned threshold at every
history length, down to 100 days per route, and approaches (on the city
routes, at 20{,}000 days, slightly exceeds) the binned high-data
programs as the history grows; those programs are estimates, not exact
optima, and we have relabelled them accordingly. This result also led us to withdraw the
recommendation in the original Remark~1 to fall back to a threshold for
histories below about 500 paths. (iii) \tref{tab:timing} reports
single-core offline fitting times for every policy, with tuning
included (median \timeBatonMs{} per route for \BATON), and the online
decision time, which is \timeOnlineUs{} per decision for every
cost-comparison policy. The RL policy is now retrained on the same CPU;
its training takes \rlTrainMin--\rlTrainMax{} minutes for the 49 routes,
against \rlBatonFitS{} seconds for fitting \BATON{} on the same routes.
\chg ``Protocol'' paragraph in \sref{sec:setup}; new
\sref{sec:budget} with \tref{tab:budget}, \tref{tab:timing} and
\fref{fig:budget}.

\begin{comment}
Statistical analysis. The experimental unit and dependence structure
should be clarified. Paired confidence intervals, effect sizes, and
route-level results would be more informative than extremely small
p-values alone. Since the problem concerns costly emergencies, tail-risk
measures would also be useful.
\end{comment}
\resp We agree, and the original pooling of 240 plan comparisons across
gates was inappropriate: the six gates re-use the same forty instances.
The revised analysis (``Statistical analysis'' in \sref{sec:setup})
takes the instance within a gate as the experimental unit. Routes of a
plan share their test days and are aggregated within the plan, and each
gate is analysed separately ($n = 40$). \tref{tab:stats} reports, per
gate and against the strongest competitor, \BATON-ho and the high-data
three-action program: the mean paired difference in saving with a 95\%
bootstrap confidence interval, the number of instances won, the
matched-pairs rank-biserial effect size, and a Holm-adjusted Wilcoxon
test. Because the strongest competitor is chosen on the same instances,
its interval re-selects the competitor in every bootstrap resample, and
its $p$-value is the largest of the $p$-values against each candidate
competitor (an intersection-union test), so the selection does not
inflate the evidence; the same selection-aware interval is used in the
robustness tables. Against the strongest competitor the advantage is
\statConsDeltaMin--\statConsDeltaMax{} percentage points on the five
conservative gates, every interval excluding zero; on the deterministic
gate \BATON{} and the equal-data three-action program are tied
($\Delta = \statDetDelta$, \statDetCI). The key contrast of each
robustness table (\tref{tab:large}, \tref{tab:dependence},
\tref{tab:pool}, and the day-type comparison in the text) now carries a
paired bootstrap interval as well, and the text reports the narrow ties
and the one loss with their intervals. Route-level results for every policy, gate
and instance are released as supplementary data. For tail risk,
\tref{tab:tail} reports the CVaR$_{95}$ of the daily plan-level bill
and the daily probability that any route of a plan suffers an
emergency; \BATON{} reduces the CVaR$_{95}$ by
\tailBatonLo--\tailBatonHi{} relative to the reactive policy.
\chg ``Statistical analysis'' paragraph in \sref{sec:setup}; new
\tref{tab:stats} and \tref{tab:tail}; the pooled $p$-values have been
removed, and the RL comparison is now tested at the instance level;
route-level CSV files in the supplementary material and the public
repository.

\begin{comment}
Methodological positioning. The manuscript should distinguish more
clearly between the novel modeling contribution, the structural
results, and the adaptation of established regression-based
optimal-stopping methods. The novelty claims should be stated
accordingly.
\end{comment}
\resp We have rewritten the contribution statement in the introduction
as three separate contributions of different kinds. The first is a
\emph{model}: execution-stage recourse in simultaneous
pickup-and-delivery as optimal stopping over a menu of levers priced per
stop from route geometry. The second is a set of \emph{structural
results} (Propositions~\ref{m-prop:bias}--\ref{m-prop:monotone}). The
third is a \emph{method} that we explicitly describe as an adaptation of
regression-based backward induction (Longstaff and Schwartz; Tsitsiklis
and Van Roy), for which we claim novelty only for the pieces the model
requires: the isotonic step licensed by
Proposition~\ref{m-prop:monotone}, the simulation-based valuation of
the post-return state, and the train-set selection of the recourse menu.
We now also cite the classical antecedents of the structural results:
monotone optimal policies for Markov decision processes (Serfozo, 1976;
Puterman, 1994) and stochastic-order comparison methods (M\"uller and
Stoyan, 2002), of which Proposition~\ref{m-prop:monotone} is an
instance; the per-stop threshold structure of optimal restocking on a
fixed route (Yang et al., 2000); and the fixed-sequence dynamic programs
for pickup and delivery with depot returns of Minis and Tatarakis (2011)
and Tatarakis and Minis (2009). We corrected the convergence claim: the
analysis of Cl\'ement et al.\ covers a fixed linear basis, and the
manuscript now points to the nonparametric analyses of Egloff (2005) and
Zanger (2013) without claiming a convergence theorem for \BATON, noting
that the regression target is monotone only in the limit because it is
computed under the fitted downstream policy. We have also toned down
several formulations (for example ``the most comprehensive computational
study'', ``significance levels that leave no room for doubt'' and
``no tuning'').
\chg Introduction, paragraphs 5--7; abstract; related work (recourse
strand); remark after the proof of Proposition~\ref{m-prop:monotone};
\sref{sec:estimation}.

\begin{comment}
Scope of the conclusions. Some managerial conclusions should be
qualified because the model does not include time windows,
heterogeneous delay penalties, shared-resource constraints, or
integrated route planning.
\end{comment}
\resp Agreed. The conclusion now contains a paragraph ``Scope of the
conclusions'' that names exactly these limits. The model has no time
windows and prices lateness uniformly per downstream customer. Routes
are fixed and executed in the planned order, and execution is not
integrated with planning. Shared resources enter only through the
standby pool of \sref{sec:pool}, and demand and prices are calibrated
rather than taken from one operator. The paragraph also explains why we
believe the \emph{comparisons} between policies are robust to these
choices, while the absolute savings are indicative. The managerial
findings themselves are now qualified where needed. For example, the
threshold critique is stated for the myopic break-even threshold, with
the conditions under which a threshold is good enough, and the value of
the depot return is attributed to plan slack and the demand profile.
We had also attributed its low value on the city instances to urban
geography; re-planning the city instances under the \textsc{SAA} gate
to test this produced plans on which breaches all but vanish (reactive
cost \citySaaReact{} per plan-day), so the two explanations cannot be
separated with our instances, and the manuscript no longer claims the
geographic one.
\chg Conclusion, new paragraph ``Scope of the conclusions'' and revised
managerial findings.

% ════════════════════════════════════════════════════════════════════════
\section*{Reviewer 2}

\begin{comment}
Independence assumption vs. experimental design. Assumption 1
(independent increments) underpins the Markov sufficiency of the scalar
load and the monotonicity proof in Proposition 3. Yet the experiments
use a Gaussian copula with equicorrelation $\rho = 0.6$. The authors
should clarify whether the backward induction and Proposition 3 remain
valid under dependent increments, or whether the method proceeds as if
independence holds. A sensitivity analysis varying $\rho$ (e.g., 0,
0.3, 0.6, 0.9) would substantially strengthen the robustness claim.
More directly: if Proposition 3 requires independence, is isotonic
regression still justified in the $\rho = 0.6$ experiments?
\end{comment}
\resp Thank you for pressing on this point. Our answer has three
parts. First, the proposition (now Proposition~\ref{m-prop:monotone})
does not require independence. Its proof needs only that the load be a
Markov chain with stochastically monotone kernels (revised
Assumption~1). This covers independent increments and the Gaussian
day-factor model, in which $W_k$ is sufficient for the factor. Second,
the benchmark copula model satisfies this only approximately, because
pickups and deliveries load on separate factors with customer-specific
scales. In that case the method proceeds as regression-based stopping
in the $(k,W_k)$ state, and the isotonic step is the monotone projection
of $\mathbb E[\,\cdot \mid W_k]$. We now say this explicitly rather than
leaving it implicit. Third, we ran exactly the sensitivity analysis you
suggest, $\rho \in \{0, 0.3, 0.6, 0.9\}$, plus a multiplicative day
factor (\tref{tab:dependence}), and tested the shape itself on fifty
thousand paths per route. The significant-violation rate of monotonicity
is \shapeViolZero{} under independence (the false-positive level of the
test), \shapeViolThree{} at $\rho = 0.3$, \shapeViolSix{} at $\rho =
0.6$, \shapeViolNine{} at $\rho = 0.9$ and \shapeViolDayfac{} under the
day factor. To show that the test is not blind, we injected a dip of
10\% and of 25\% of the mean continuation cost into one bin; it is
detected in only \shapePowTenAll{} and \shapePowTwentyFiveAll{} of the
tests. The test therefore finds no significant violations, but with
limited power: small dips could go unnoticed. We read the result as
consistent with a monotone continuation value, and the isotonic
restriction as a reasonable approximation where the theorem does not
strictly apply, rather than as a proof of monotonicity. A \BATON{} variant keeps the highest saving
among implementable policies under all five laws. \BATON{} itself is
tied with its strongest competitor in three of the ten law--gate rows and
loses one, at $\rho = 0.9$ on \textsc{Det} plans, to the equal-data
three-action program (\depRhoNineDetDelta{} points,
\depRhoNineDetDeltaCI); the state-conditional variant reverses that
loss. On the \textsc{Det} plans \BATON's share of the high-data
three-action saving is \depRhoZeroDetRatio, \depRhoSixDetRatio{} and
\depRhoNineDetRatio{} at $\rho = 0$, $0.6$ and $0.9$, and under
independence it attains \exDetShare{} of the \emph{exact} optimum,
which we now compute by numerical convolution (\tref{tab:exact}).
Interestingly, all implementable policies save \emph{more} as $\rho$ grows:
correlation makes the running load informative about the rest of the
day.
\chg Revised Assumption~1 and discussion (\sref{sec:formulation}); proof
of Proposition~\ref{m-prop:monotone} and the paragraph after it
(\sref{sec:estimation}); \sref{sec:robust} with \tref{tab:dependence}
and \tref{tab:exact}.

\begin{comment}
Bias in the fresh-start value estimator. In \S3.5, $F_k = C_k(0)$ is
estimated by simulating the reset state forward through the
already-fitted downstream policy, not the optimal policy. Since the
fitted policy is weakly suboptimal, the estimator is an upper bound on
the true $F_k$, which systematically overprices $R_k + F_k$ and biases
the policy away from depot return---potentially toward handoff or
continuation. The authors should discuss the direction and magnitude of
this bias, and quantify it in settings where depot return should be
valuable (e.g., Salhi--Nagy CMT geometry).
\end{comment}
\resp This is an excellent question, and answering it led to an
improvement of the method. The reviewer is right that the suboptimality
of the downstream policy pushes $\widehat F_k$ upward. In the revision
we identify two further sources of bias, of opposite sign, and quantify
all three against a high-data reference fresh-start value computed from fifty
thousand independent paths (\sref{sec:actions} and \tref{tab:fresh}).
(i) The suboptimality you describe pushes $\widehat F_k$ up. We note
that this is also the value the deployed policy will actually realise
after a return, so it is the consistent price for the policy that is
deployed. (ii) The forward pass runs on the training paths on which the
downstream models were fitted, which pushes $\widehat F_k$ down by
in-sample optimism. (iii) Under positively dependent demand, returns are
chosen on heavy days whose remainder is also heavy, so an unconditional
$F_k$ under-prices the return on exactly the days on which it is used.
On Salhi--Nagy, the geometry you suggest, the in-sample error of
$\widehat F_k$ against the reference value is \frSalhiBiasIn{} of the
handoff price and the out-of-sample error of the same fitted policy is
\frSalhiBiasOut. Replacing $\widehat F_k$ by the reference value
changes the three-action saving from \frSalhiThree{} to
\frSalhiFstar. Sources (i) and (ii) are thus small and nearly cancel.
Source (iii) is the material one, and it appears on delivery-dominated
city routes, where returns rarely pay. It has a direct remedy within
the method: regressing each training path's realised post-return cost
on its own pre-return state gives a state-conditional fresh-start value
$\widehat F_k(W_k)$, monotone for the same reason as $C_k$ and flat
under independence. The resulting \textsc{Baton-cf} needs no deployment
selection and improves on \BATON{} on all six Dethloff gates
(\tref{tab:grand}) and in \depCfAbove{} of the ten rows of
\tref{tab:dependence}. Since it was introduced during the revision, after
the city results had revealed bias (iii), we report it beside \BATON{} in
the main tables, say so in the text, and keep \BATON{} as the headline
policy so that the main tables remain comparable with the submitted
version. The high-data three-action reference DP$^3_{50\mathrm{k}}$ uses
the same state-conditional fresh-start value, which the manuscript now
documents.
\chg New paragraphs ``Estimating the fresh-start value'' and ``Deployment
selection'' in \sref{sec:actions}; \sref{sec:robust} and
\tref{tab:fresh}; \textsc{Baton-cf} in \tref{tab:grand},
\tref{tab:large}, \tref{tab:dependence} and \tref{tab:timing}.

\begin{comment}
Action set and the value of the third lever. The paper only considers
two levers: full handoff of the remaining route and depot return.
Partial handoff---transferring only the riskiest remaining
customers---is mentioned as future work but is likely common in
practice. The claim that the action set is ``general in form'' needs
clearer boundaries. Moreover, Table 4 shows that in real city instances
BATON performs essentially the same as its handoff-only restriction
(10.8\% vs. 10.7\%), implying that depot return is rarely useful on real
road networks. The authors should discuss this more explicitly: the
three-action policy's value appears mainly in synthetic benchmarks and
conservative planning gates, not in real urban settings.
\end{comment}
\resp We agree on both counts. On the boundaries: Proposition~\ref{m-prop:monotone}
and the algorithm cover exactly the actions that either end the route
at a price known at the current stop, or reset the load to a known
level at a known price. The handoff and the depot return are of these
kinds; a partial handoff is not, because it changes the remaining route
itself. The conclusion now states this. It also explains what a partial
handoff would require: a state augmented by the split point, a price
schedule for every split, and one regression per stop and split, at a
cost quadratic rather than linear in the route length. On the value of
the third lever: we now say plainly, in \sref{sec:large} and in the
conclusion, that the depot return is valuable where plans carry slack
and depots are central---the conservative Dethloff gates and the
Salhi--Nagy geometry, where \BATON{} saves \svSalhiBaton{} against
\svSalhiHo{} for \BATON-ho---and that on our city instances it rarely
pays. There \BATON{} keeps the return in the menu on only
\cityActShare{} of routes and equals its handoff-only form
(\svCityBaton{} against \svCityHo). We had attributed this to urban
geography. To test that explanation we re-planned the city instances
under the \textsc{SAA} gate, which leaves slack; but on those plans
breaches all but vanish (reactive cost \citySaaReact{} per plan-day), so
no policy has anything to save. The city results are equally consistent
with the demand profile (delivery-dominated routes whose load drifts
downward) and the lack of slack in \textsc{Det} plans---on the Dethloff
\textsc{Det} plans the return also adds only \detActGain{} points. The
manuscript now states the finding in these terms and no longer claims
that urban geography as such makes the return worthless.
\chg Conclusion (managerial finding four; ``Limitations and next
steps''); \sref{sec:large}.

\begin{comment}
Quantitative characterization of Proposition 2. The proof shows that a
fixed threshold over-triggers, but does not quantify by how much. The
authors should discuss when the threshold family remains ``good
enough'' and whether the cost of over-triggering can be small in
practice. This would help readers judge whether the structural critique
is mainly qualitative or also economically significant.
\end{comment}
\resp We thank the reviewer for this suggestion, which led to a new
result. The new Proposition~\ref{m-prop:regret} shows that the myopic
stopping time never exceeds the optimal one, pathwise. The excess
expected cost of the myopic rule equals
$\mathbb E[(H_{\sigma^0} - C_{\sigma^0}(W_{\sigma^0}))
\mathbf 1\{\sigma^0 < \sigma^\star\}]$, and is bounded by
$\mathbb E[(H_{\sigma^0} - L_{\sigma^0})\mathbf 1\{\sigma^0 <
\sigma^\star\}]$, where $L$ is the clairvoyant cost of the rest of the
day. Under flat prices this is at most the handoff price times the
probability of an intervention on a day that would have completed
cleanly. The threshold is therefore good enough exactly when the states
in which it fires almost surely lead to a breach and prices decline
slowly---tightly packed plans and short routes---and poorer on plans
that carry slack. We also corrected the scope of the critique: the
proposition concerns the myopic break-even threshold, and the abstract
and introduction no longer claim that every fixed threshold
over-triggers. \tref{tab:regret} evaluates both sides on our routes.
The pathwise ordering holds on \regretContained{} of test days and the
bound on \regretBoundHolds{} routes; on the remaining route the regret
exceeds the bound by \regretFailMax{} of reactive cost, within Monte
Carlo error. Pooled over routes, the myopic regret is \regretAll{} of
reactive cost against a bound of \regretBoundAll{} (the average of the
route-level ratios is \regretRouteMean), and the clean-day term is
\regretCleanShare{} of the bound, so unnecessary interventions are
indeed what is paid for. The regret is smallest on the city plans
(\regretCity) and on short routes (\regretShort{} for at most 12 stops),
\regretDet{} on Dethloff \textsc{Det} plans, and larger on conservative
\textsc{SAA} plans (\regretSaa) and longer routes. A threshold tuned on
realised cost is a different object---tuning moves it away from the
break-even level, so the proposition does not bound it---but its excess
over \BATON-ho, fitted on the same thousand days, is small on tight
plans too: \thrGapDet{} of reactive cost on \textsc{Det} plans and
\thrGapCity{} on the city plans, against \thrGapSaa{} on \textsc{SAA}
plans. The answer to the reviewer's question is therefore: for the
handoff decision alone, the cost of over-triggering is economically
modest on tight plans---there a tuned threshold is good enough---and
grows on plans with slack. The large differences in \tref{tab:grand} on
the conservative gates come mainly from the depot return, which a
handoff threshold cannot price, and a two-lever threshold rule that is
given the return recovers only part of them.
\chg New Proposition~\ref{m-prop:regret} and discussion
(\sref{sec:defects}); ``The price of over-triggering'' in
\sref{sec:learning} with \tref{tab:regret}; managerial finding two in the
conclusion.

\begin{comment}
Planning--execution co-optimization. [\ldots] Since this
co-optimization is only an outlook, the authors should state whether
they consider it the most important future direction and why it was not
implemented here.
\end{comment}
\resp Yes, we consider it the most important direction, because it
attacks the breach risk that no execution policy can reach: the risk at
the first stop of a packed route. We did not implement it for two
reasons, which the conclusion now states. First, it changes the object
of study. Plans would depend on the execution policy, so the design
that isolates this paper's contribution---comparing execution policies
on common plans---would no longer be available. Second, it requires the
fitted execution cost of every candidate route inside the planner's
inner loop, which is a separate computational question.
\chg Conclusion, last paragraph.

\begin{comment}
Demand exchangeability. The model treats days as exchangeable and does
not consider seasonality, day-of-week effects, or promotional days.
[\ldots] A brief discussion would be useful.
\end{comment}
\resp We added both a discussion and an experiment (\tref{tab:daytype}).
One day in five is a promotion day with higher deliveries and
considerably higher pickups, and we compare three strategies on the
same history of $10^3$ days: a pooled fit that ignores the day type;
day-type-specific fits that split that history by type (about
\dtNNormal{} normal and \dtNPromo{} promotion days), as an operator can
when promotions are scheduled; and a stale fit that uses the normal days
only and is applied unchanged on promotion days. On promotion days
\BATON{} saves \dtPromoAware{} with the day-type-specific fit,
\dtPromoPooled{} with the pooled fit and \dtPromoStale{} with the stale
fit. Over all days the day-type-specific fits are no better than the
pooled fit (\dtAllAware{} against \dtAllPooled; per plan
\dtAwarePooledDelta{} points, \dtAwarePooledCI): what they gain on
promotion days they lose on normal days by fitting on fewer paths. The
pooled policy copes because a promotion day reveals itself through a
heavier running load, which the monotone continuation value prices
whatever the day type. The conclusion draws the practical rule: keep
known day types in one training history; fit separately only when each
type has a long history of its own; and where the mix drifts, re-fit on
a rolling window and monitor.
\chg \sref{sec:robust} with \tref{tab:daytype}; conclusion,
``Limitations and next steps''.

\begin{comment}
RL comparison. [\ldots] The authors may briefly discuss whether other RL
architectures (e.g., DQN, PPO, attention-based) might narrow the gap,
without necessarily running additional experiments.
\end{comment}
\resp We added this discussion to \sref{sec:learning}. Value-based
methods such as DQN estimate the same continuation value with a generic
function class, and we would expect them to approach, not exceed, the
isotonic regression at equal data. Actor-critic methods such as PPO
would likely close part of the gap by learning a critic that plays the
role of $\widehat C_k$. Attention- and graph-based architectures pay off
when the state is a large combinatorial object, and would become
relevant for partial handoffs or fleet-level coordination. We also
state in the manuscript that our three seeds and two epoch budgets do not
rule out that longer training, other hyper-parameters or another
architecture would narrow the gap; sample efficiency is our
interpretation for the architecture we tested. Finally, we
retrained the RL baseline on a route bundle consistent with the current
plans, on the same CPU as \BATON, and report its training time.
\chg \sref{sec:learning}, paragraph after \tref{tab:rl}; related-work
paragraph on learning.

\begin{comment}
Clairvoyant bound. BATON exceeds the handoff-only clairvoyant bound on
four of six planning gates. Section 4.2 explains this correctly, but the
explanation should appear more prominently in the abstract and
conclusion to avoid confusing readers.
\end{comment}
\resp Done. The abstract now says that, because the depot return
enlarges the action set, \BATON{} can exceed the clairvoyant bound of the
handoff-only class. The explanation---the bound is clairvoyant only
about the handoff lever, and a return can prevent a breach at a price no
handoff can match---now appears in the introduction, in
\sref{sec:algorithms} where the bound is defined, in \sref{sec:headline},
and in the first paragraph of the conclusion. Every results table
labels the oracle as handoff-only (``oracle-ho''). We have also added
the clairvoyant bound for the full three-action menu (oracle$^3$), which
lies above every policy; \BATON{} attains
\ratioOrcThreeLo--\ratioOrcThreeHi{} of its saving on the Dethloff
gates.
\chg Abstract; introduction; \sref{sec:algorithms}; \sref{sec:headline};
conclusion; table captions.

\begin{comment}
Zero-pickup stations. City instances use pickup fractions uniform on
[0.2, 0.8], excluding stations with zero pickup. Since deliver-only
customers are common in practice, the authors should either justify
this choice or add a sensitivity check with a nonzero fraction of
zero-pickup stations.
\end{comment}
\resp We added the sensitivity check. Every city instance now has two
deliver-only twins in which the pickup of a random 25\% or 50\% of the
customers is set to zero, and the twins are re-planned
(\tref{tab:large}). With fewer pickups the load rises less along the
route, breaches are rarer, and all savings shrink: at 25\% and 50\%
deliver-only customers \BATON{} saves \svZpTwentyFiveBaton{} and
\svZpFiftyBaton{}, while the tuned threshold falls to \svZpFiftyThr{} at
50\%. At 25\% \BATON{} still leads its strongest competitor by
\lgZpTwentyFiveDelta{} points (\lgZpTwentyFiveDeltaCI); at 50\%, where
recourse is worth least, it is tied with the equal-data dynamic program
(\lgZpFiftyDelta{} points, \lgZpFiftyDeltaCI) and reaches only
\lgZpFiftyRatio{} of the high-data reference---the lowest share in the
study, which we report in the abstract's range. Our original choice of
strictly positive pickups followed the construction of the Dethloff
benchmark, in which every customer has both a delivery and a pickup; the
new twins show that the ranking survives deliver-only customers while
the size of the prize, and \BATON's margin, shrink.
\chg \sref{sec:setup} (instances); \tref{tab:large}; \sref{sec:large}.

\begin{comment}
Figure 2. The three panels are small and may be difficult to read in
print. Consider splitting into two figures or increasing font size.
\end{comment}
\resp We did both: the explainer is now split into two figures
(\fref{fig:explainer}, two panels at full text width, and
\fref{fig:explainer2}, the test-day costs), with all fonts enlarged to
print size.
\chg \fref{fig:explainer} and \fref{fig:explainer2}.

% ════════════════════════════════════════════════════════════════════════
\section*{Reviewer 3}

\begin{comment}
1. The paper divides the vehicles into three classes, the planned
fleet, the standby class and the emergency class. What is the
difference between the standby and the emergency vehicles? How many
standby vehicles are there?
\end{comment}
\resp The two classes differ in how they are contracted and in what
they cost the customer, not in the vehicles themselves. A standby
vehicle is reserved capacity under a pre-agreed rate: it is dispatched
before a breach, takes over the remaining route at a small surcharge,
and keeps the downstream customers on schedule. An emergency vehicle is
hired on the spot market after the planned vehicle has overflowed: it
pays surge rates, must first reach a stranded vehicle, and every
downstream customer is served late. In the default model the standby
pool is shared across the operation and billed per vehicle-day used, so
the number of standby vehicles is not a constraint but a consequence
of the policy. We now report it: the pool that covers 95\% of days has
a median of \poolDetMed{} and a maximum of \poolDetMax{} vehicles per
plan for \BATON{} on the Dethloff \textsc{Det} plans, and a median of
\poolSAAMed{} on the \textsc{SAA} plans, where the depot return does
most of the work (over all six gates: median \poolBatonMed, maximum
\poolBatonMax; \poolHoMed{} and \poolHoMax{} for the handoff-only
restriction). In
the new experiment with a pool of fixed size reserved in advance
(\sref{sec:pool}), reserving vehicles pays only when they can be held
for well under the standby day rate and handoffs are frequent: at a
holding cost of 5 per vehicle-day the optimal pool is about
\poolDethSstarLow{} vehicles for a packed Dethloff \textsc{Det} plan
and about \poolMetroSstarLow{} for a metropolitan pool serving some
\poolMetroK{} routes, while slack plans need none.
\chg New paragraph ``Standby versus emergency vehicles'' in
\sref{sec:setting}; \sref{sec:pool} and \tref{tab:pool}.

\begin{comment}
2. Since the paper does not consider the holding cost of the standby
vehicles, if the number of the standby vehicles larger than the
emergency vehicles, then the Assumption 2 and the Eq. 7 may not hold.
\end{comment}
\resp Thank you; the original text did not make this clear. In the
default model the holding cost enters through the standby day rate
$F_{\mathrm{sb}}$ in the handoff price, charged per vehicle-day
consumed. The revised
\sref{sec:setting} says so explicitly. When standby vehicles are
instead reserved in advance, their holding cost is paid whether or not
they are used; the new experiment of \sref{sec:pool} models exactly
this, with $S$ reserved vehicles, a range of holding costs and the
break-even holding cost of a reserved vehicle.

The reviewer is right that Eq.~(7) ($E_k > H_k$) can fail if standby
capacity is expensive. We have therefore checked what the theory
actually needs, and it is less than we had assumed. The revised
Assumption~2 requires only that the emergency schedule be
non-increasing along the route; no ordering between $H_k$ and $E_k$ is
used anywhere. The proof of Proposition~\ref{m-prop:monotone} bounds
the value function by the emergency price through the cost of never
acting, not through $H_k$. Eq.~(7) is now presented as a property of
our calibration, not a requirement. When it fails, the handoff is simply
dominated at the stops concerned and the policy never uses it there.
To demonstrate this we added a configuration with $F_{\mathrm{sb}} = 60$,
in which the handoff price exceeds the emergency price at late stops
(\tref{tab:costsens}, the $F_{\mathrm{sb}} = 60$ row). \BATON{} saves \sbSixtyBaton{} there,
against \sbSixtyThresh{} for the tuned threshold.
\chg Revised Assumption~2 and the paragraph after it; text after
Eq.~\eqref{m-eq:prices}; proof of Proposition~\ref{m-prop:monotone};
new row of \tref{tab:costsens}.

\begin{comment}
3. In Sec 4.2, it is claimed that `BATON attains the lowest execution
cost on every gate', but in Table 2 and Fig. 4, it is indicated that in
the case of `Deterministic' and `W-DRO', the value of `oracle' is
greater than that of `BATON'.
\end{comment}
\resp Thank you; our wording was imprecise. The oracle is a clairvoyant
policy: it knows the whole day's demand in advance, so it cannot be
implemented, and it serves only as a bound. The claim should have read
``the lowest execution cost of every \emph{implementable} policy''. It
now does, everywhere (abstract, introduction, \sref{sec:headline},
conclusion), and the main tables now separate the reference points (the
clairvoyant bounds and the high-data dynamic programs) from the
competitors. The
oracle is moreover restricted to the handoff lever, which is why
\BATON{} exceeds it on the other gates: a depot return can prevent a
breach at a price no handoff can match. On the \textsc{Det} and
\textsc{WDRO} gates this second lever is worth less, and the
clairvoyant handoff-only policy stays ahead.
\chg \sref{sec:headline} (first and second observations); captions of
\tref{tab:grand} and \fref{fig:gatedots}; abstract and conclusion.

\begin{comment}
4. In Table 3 and Table 5, only the value of `BATON-HO' is displayed.
What is the value of `BATON'?
\end{comment}
\resp Both tables now show \BATON{} as well. In the synthetic scenarios
(\tref{tab:synthetic}) the routes have no geometry from which to price a
depot return; the submitted version therefore used the handoff lever
only. We now price the return flat at half a handoff and report the full
\BATON, which reaches \ctdBatonFull{} on the collect-then-deliver
scenario. In the RL comparison (\tref{tab:rl}) the learned policy acts
on the handoff lever only, so \BATON-ho is the like-for-like comparison
(\rlBatonHo); the full \BATON, which adds the depot return, saves
\rlBaton{} on the same routes.
\chg \tref{tab:synthetic} and \tref{tab:rl}.

\begin{comment}
5. Table 4 shows that in the case of `City, real shops' and `City,
uniform', the value of `BATON' is smaller than that of `DP50k' and
`oracle'. Why?
\end{comment}
\resp Thank you for pointing this out; the original tables did not make
the role of these two columns clear. Both are reference points rather
than competitors, and the revised tables separate them accordingly. The
\emph{oracle} knows the day's demands in advance. On city routes much of
the breach risk comes from demand surprises at the next stop that no
policy observing only the past can anticipate, so its lead is large and
unattainable. DP$_{50\mathrm{k}}$ is fitted on fifty times more data
than \BATON. On the city instances the depot return rarely pays, so the
comparison is between two handoff-only policies, and the gap
(\svCityBaton{} against \svCityDp) is a gap to a reference that has
fifty times more data and also a different estimator (binned means
instead of isotonic regression), so it cannot be attributed to the data
alone. The new data-budget experiment indicates that the amount of data
accounts for most of it (\tref{tab:budget}, \fref{fig:budget}): given
20{,}000 training days, \BATON{} saves \budCityLarge{} on the city
instances, slightly more than DP$_{50\mathrm{k}}$ (\budCityDp; the
binned program is an estimate, not an exact optimum). The oracle's lead, by contrast, does not shrink with
data: it is the value of knowing the future.
\chg \sref{sec:large}; \sref{sec:budget} with \tref{tab:budget} and
\fref{fig:budget}; table captions.

% ════════════════════════════════════════════════════════════════════════
\section*{Corrections found by the authors}

\begin{enumerate}[leftmargin=1.4em]
\item \textbf{Proposition 1 (endpoint-label bias).} The proof of the
  cost bound stated that the optimal policy is ``no worse than the
  clairvoyant policy''; the inequality runs the other way, so the
  claimed upper bound on the optimal cost and the ``factor
  $C_{\mathrm{fail}}/\omega_F$'' were not proved. The corrected
  Proposition~\ref{m-prop:bias} bounds the loss of the endpoint-trained
  policy from both sides, with the upper bound attained when the breach
  is announced one stop ahead.
\item \textbf{Proposition 2 (over-triggering).} The strictness condition
  in the original statement was not sufficient: the policy used in the
  proof also hands off on days that would not have breached. The
  corrected Proposition~\ref{m-prop:myopic} gives a condition that is
  proved---the path must reach, with positive probability, a later
  state in which the myopic rule itself fires. It also derives strict
  inclusion of the stopping regions under continuity, and it is now
  stated for position-dependent prices.
\item \textbf{A side condition in Proposition 1.} The formal
  verification showed that the lower bound on the loss needs routes with
  at least two stops (so that a decision epoch precedes the last stop)
  and $0 \le \omega_F \le C_{\mathrm{fail}}$; the statement now says so.
\item \textbf{The ``92--99\% of a near-exact dynamic program'' claim}
  in the original abstract did not hold on every gate: it referred to
  the handoff-only program, and on the \textsc{WDRO} and
  \textsc{M-DRO} gates the ratio was lower. Moreover the binned
  programs are estimates, not exact optima. The revised abstract compares
  \BATON{} with the exact optimum on deterministic plans under independent
  demand (\exDetShare) and gives the full range against the high-data
  reference over all benchmarks and demand laws with a non-negligible
  reactive cost (\ratioAllLo--\ratioAllHi); the text adds the share on
  the Dethloff gates (\ratioDpThreeLo--\ratioDpThreeHi) and the exact
  share on slack plans (\exSaaShare).
\item \textbf{``Every fixed threshold over-triggers.''} The original
  abstract, introduction and conclusion attributed to every fixed
  threshold what Proposition~\ref{m-prop:myopic} proves for the myopic
  break-even rule. The claims are now stated for that rule; a threshold
  tuned on realised cost is not covered.
\item \textbf{Convergence citation.} The analysis of Cl\'ement et al.\
  (2002), which we cited for the convergence of the backward induction,
  covers a fixed linear basis and not isotonic regression; the text now
  cites the nonparametric analyses of Egloff (2005) and Zanger (2013)
  and states that we do not prove a convergence theorem.
\item \textbf{Urban geography.} The original text attributed the low
  value of the depot return on the city instances to congested urban
  networks. Our instances cannot separate geography from plan slack and
  demand profile (see the reply to Reviewer~2, comment~3), and the
  claim has been withdrawn.
\item \textbf{Reproducible scenario seeds.} The evaluation derived its
  per-instance seeds from Python's salted string hash, so scenario draws
  differed from run to run. Seeds are now fixed functions of the
  instance name, and all results were regenerated. Numbers moved by
  sampling noise (for example, the tuned threshold's saving on
  \textsc{WDRO} plans is now \svWDROThresh{} rather than $-8.3\%$); no
  qualitative conclusion changed.
\item \textbf{Cost-sensitivity baseline.} The baseline row of the
  sensitivity table was computed on all 40 instances while the sweep
  used 12; it is now computed on the same 12.
\item \textbf{RL baseline.} The route bundle of the RL experiment had
  been built from an earlier version of the city plans. It is rebuilt
  from the current plans (49 routes), and the RL policy is retrained on
  it. The description of the training now matches the code; the
  original text mentioned entropy regularisation and a learning-rate
  search, which the code did not perform.
\item \textbf{Timing.} The original text said ``single-digit
  milliseconds'' for fitting; the measured median for the full policy
  with deployment selection is \timeBatonMs{} per route, which we now
  report.
\item \textbf{W-DRO gate.} With a type-1 Wasserstein ball and a
  1-Lipschitz loss, the worst-case CVaR equals the empirical CVaR plus
  $\varepsilon/(1-\alpha)$, so the \textsc{WDRO} gate is the
  \textsc{SAA} gate at a tighter threshold. We now say so in
  \sref{sec:setup} and present it as a more conservative regime rather
  than an independent paradigm.
\end{enumerate}

We hope that the revised manuscript meets the reviewers' expectations,
and we thank them again for their time and insight.

\medskip
\noindent Sincerely,\\
Quang-Vinh Dang (corresponding author), on behalf of all authors

\end{document}
